% Shared preamble for the English summaries in this repository.
\documentclass[11pt,a4paper]{article}
\usepackage[margin=2.2cm]{geometry}
\usepackage[T1]{fontenc}
\usepackage[utf8]{inputenc}
\usepackage{lmodern}
\usepackage{amsmath,amssymb}
\usepackage{graphicx}
\usepackage{booktabs}
\usepackage{enumitem}
\usepackage[hidelinks]{hyperref}
\setlist{itemsep=2pt,topsep=4pt}
\setlength{\parindent}{0pt}
\setlength{\parskip}{5pt}
\graphicspath{{figures/}}
\newcommand{\summaryheader}[3]{%
  \begin{center}
    {\Large\bfseries #2}\\[3pt]
    {\large #1}\\[3pt]
    {\small Arshia Goshtasbi \textbullet{} MSc Electrical Engineering (Control), Iran University of Science and Technology \textbullet{} #3}
  \end{center}
  \vspace{2pt}\hrule\vspace{8pt}
  \noindent\textit{English summary of my original Persian report, which is kept as a PDF in this folder.}\par}

\begin{document}
\summaryheader{Convex Optimization, Homework 3}{Convex and Quasi-convex Functions}{Spring 1404 (2025)}

This homework is about deciding whether functions are convex, concave or quasi-convex, using second-order conditions and the composition rules.

\section*{Problems and what I did}
\begin{enumerate}
  \item \textbf{Running average of a function.} I showed by differentiation that $g(x)=\tfrac{1}{x}\int_0^x f(t)\,dt$ inherits monotonicity and concavity from $f$.
  \item \textbf{Epigraphs.} I matched the epigraph types to the functions: a half-space comes from an affine function, a convex cone from a positively homogeneous function, and a polyhedron from a piecewise-linear function.
  \item \textbf{Hessian tests.} For each function I computed the Hessian. When it was indefinite, I checked the sub- and superlevel sets instead: some functions are quasi-concave only, some quasi-linear (both quasi-convex and quasi-concave), and some convex and concave as well as quasi-convex.
  \item \textbf{A function of several variables.} I computed the Hessian entries for $i\neq j$ and $i=j$, and used an inequality between two auxiliary vectors $a$ and $b$ to show that the function is concave.
  \item \textbf{Composition rules.} I used them for
    \begin{itemize}
      \item the sum of the $k$ largest residuals of $Ax-b$, a pointwise maximum of convex functions;
      \item $-\log(-g(x))$ for convex $g<0$, where $-\log$ is convex and decreasing;
      \item $y^Ty/t$ composed with an affine map, a perspective of a quadratic;
      \item a $\log$ of a sum of exponentials, through Jensen's inequality.
    \end{itemize}
  \item \textbf{Products and ratios.} Using Jensen's inequality, I showed that the product of two non-negative convex functions that are both non-decreasing (or both non-increasing) is convex, derived the matching results for concave functions, and handled $f/g$ through $1/g$.
  \item \textbf{More compositions.}
    \begin{itemize}
      \item $\|Ax-b\|$ is convex as a norm of an affine map.
      \item $\|Ax+b\|/(c^Tx+d)$ is convex on $c^Tx+d>0$ as the perspective of a norm.
      \item $\max\{1,x,x^2\}$ is convex as a pointwise maximum.
      \item A function whose second derivative changes sign at $x=-2$ is convex for $x>-2$ and concave for $x<-2$.
    \end{itemize}
  \item \textbf{Perspective-based functions.}
    \begin{itemize}
      \item $t - \tfrac{1}{t}x^Tx$ is concave.
      \item $-\log(t - x^Tx/t)$ is convex.
      \item A sum of two negated logarithms of concave functions is convex.
    \end{itemize}
  \item \textbf{Linear-fractional functions.} I showed that $u(x)/v(x)$ with affine $u$ and positive affine $v$ is quasi-convex, since each sublevel set is a half-space. Whether it is actually convex depends on the sign of a remaining term.
  \item \textbf{A three-variable function.} $(x-z)^2/(y+1) + \max\{1+2x-y,\ 1/z,\ x^2\}$ is convex on $y>-1$, $z>0$. The first term is a quadratic-over-linear, and the second is a pointwise maximum of convex functions.
\end{enumerate}
\end{document}
